\section{Infinite Ramsey Theory}

% [CORRECTED] was: ``indeed, if I have a red/blue colouring of $\N$, then either the set of red numbers or the set of blue numbers must be infinite, so just look at the edges between them to get the desired clique.'' --- that colours the vertices rather than the edges, and an infinite set of vertices says nothing about the colours of the edges among them (colour $ij$ by whether $i$ and $j$ have the same parity: every vertex sees both colours infinitely often, and the monochromatic infinite cliques are the evens and the odds); the argument below is the standard one
Consider the complete graph on the natural numbers, denoted $K_{\N}$. First, observe that any $2$-colouring of the edges of $K_{\N}$ has an infinite monochromatic clique: indeed, pick $v_1 \in \N$; infinitely many of the edges at $v_1$ have the same colour, say $c_1$, so let $S_1$ be the infinite set of their other endpoints. Pick $v_2 \in S_1$; infinitely many of the edges from $v_2$ into $S_1$ have the same colour, say $c_2$, and let $S_2 \ssq S_1$ be the set of their other endpoints. Continuing in this way gives $v_1, v_2, \ldots$ and colours $c_1, c_2, \ldots$ such that the edge $v_i v_j$ has colour $c_i$ whenever $i < j$. One colour occurs infinitely often among the $c_i$, and the $v_i$ with that $c_i$ form an infinite monochromatic clique.

Ramsey numbers of finite sets have a strengthening in which the monochromatic set is asked to be large relative to its own elements.

% The definition below was rewritten on the author's instruction (``what the heck is this definition''). It read: Define $\bar{R_k^{(c)}}$ to be the smallest $N \in \N$ such that any $c$-colouring of $k$-sets of $[N]$ has a monochromatic $n$-clique $X$ on $n' \geq n$ such that $n' \geq \min(X)$.
\begin{boxdefinition}[Large Ramsey Number]
    A finite set $X \ssq \N$ is \textbf{relatively large} if $\abs{X} \geq \min X$. For $n, k, c \in \N$, define $\bar{R}^{(c)}_k(n)$ to be the smallest $N \in \N$ such that any $c$-colouring of the $k$-subsets of $[N]$ has a relatively large set $X \ssq [N]$ with $\abs{X} \geq n$ all of whose $k$-subsets have the same colour.
\end{boxdefinition}

So $X$ is a monochromatic clique in the sense of hypergraph Ramsey numbers, and for $c = 2$ the definition without the words ``relatively large'' would be $R_{(k)}(n, n)$. The extra condition is what makes the number interesting: a monochromatic clique whose least element is $m$ now has to have at least $m$ elements, so cliques that start early are cheap and cliques that start late are expensive. That these numbers exist at all follows from the infinite version of Ramsey's theorem by a compactness argument, which we will not give.

One interesting theorem (that we will not prove here) is that the existence of the above cannot be proven in Peano arithmetic. This is the Paris--Harrington theorem~\cite{ParisHarrington}, the first example of a natural combinatorial statement that is true but not provable from Peano's axioms.
